\begin{abstract}
An adversary watching a power-grid control network can tell which device it is looking at from
how long that device takes to answer, because vendors implement the same function differently.
Hiding that timing is challenging. The relay cannot be modified, and existing traffic obfuscation
was built for a different feature: it reshapes flow-level sizes and inter-packet times, whereas
the feature that identifies the device here is one interval inside a single transaction. It also
hides the packets it adds behind encryption, which these plaintext networks do not have.

To change the timing without touching either endpoint, we put the defense in the network. A
programmable switch in front of the outstation holds the device's acknowledgment and its
response, and releases each at a deadline computed from the same instant, the moment the outstation
acknowledges the request. Neither release depends
on how long the device took, so the interval an observer measures is a value the operator sets.
The switch forwards the packets it received, so both endpoints run unmodified and no DNP3 byte
changes.

We implement this on an Intel Tofino-1 and evaluate it against one physical SEL-751A protective
relay. The measured interval settles on the configured value: its median is the configured 4~ms,
its interquartile range falls by more than two orders of magnitude, and its sample variance falls
from 6.808 to 0.394~ms$^2$. A classifier trained on that interval before deployment is left at
chance across the three transaction classes. An adversary who retrains recovers much of that from
a second interval we do not target, and that residual bounds what the framework achieves. The
evaluation covers transaction-class timing for a single device; we do not measure confusion
between device models, and we do not claim the feature is removed for an adversary who accumulates many
observations.
\end{abstract}
